%%%%%%%%%%  scan idx 378  |  folio 365  %%%%%%%%%%
\SGAPrintPage{365}%
\sgaStableTarget{sga6.e06.expose}\section*{\underline{Exposé VI}\\[3pt] \underline{\underline{THE \(K^{\bullet}\) OF A PROJECTIVE BUNDLE:}}\\[3pt] \underline{\underline{COMPUTATIONS AND CONSEQUENCES}}}
% Exposé VI: reviewed local symmetric-algebra and Proj notation.
\renewcommand{\Sym}{\mathbf{S}}
\renewcommand{\Projz}{\underline{\operatorname{Proj}}_0}
\begin{center}
\underline{by Pierre BERTHELOT}
\end{center}

\sgaStableTarget{sga6.e06.sec.1}\section*{1. Computation of the \(K^{\bullet}\) of a projective bundle: the case of locally free sheaves of finite type}

In this section, the notation \(K^{\bullet}(X)\) will denote the Grothendieck group of locally free sheaves of finite type on a scheme \(X\).

The aim of this section is to prove the following theorem:

\sgaStableTarget{sga6.e06.thm.1-1}\textbf{Theorem 1.1.} \uline{Let} \(S\) \uline{be a quasi-compact scheme}, \(\underline{E}\) \uline{a} \(\underline{O}_S\)\uline{-Module locally free of rank} \(r+1\), \(X=\mathbf P(\underline{E})\) \uline{the associated projective bundle}; \uline{the structural morphism}: \(X\xrightarrow{f}S\) \uline{defines a ring homomorphism}: \(K^{\bullet}(S) \xrightarrow{f^*} K^{\bullet}(X)\) \uline{endowing} \(K^{\bullet}(X)\) \uline{with the structure of a} \(K^{\bullet}(S)\)-\uline{algebra}. \uline{Let} \(\xi\in K^{\bullet}(X)\) \uline{be the class of the sheaf} \(\underline{O}_X(1)\). \uline{Then} \(K^{\bullet}(X)\) \uline{is a free} \(K^{\bullet}(S)\)-\uline{module with basis} \(1,\ldots,\xi^r\).

\uline{Plan of the proof}:

\(1^\circ\)) The powers of \(\xi\) generate \(K^{\bullet}(X)\) as a \(K^{\bullet}(S)\)-module.

To prove this, we shall establish the following points:

\begin{enumerate}[label=\alph*)]
\item If \(\underline{F}\) is a \(\underline{O}_X\)-Module locally free of finite type, there exists a graded \(\Sym(\underline{E})\)-Module \(\underline{M}\) of finite presentation, flat over \(\underline{O}_S\), such that \(\underline{F}=\Projz(\underline{M})\) (\(\Sym(\underline{E})\) being the symmetric algebra of \(\underline{E}\); for the notation \(\Projz\), see EGA II 8.12.1).

\item There exists a finite graded resolution of \(\underline{M}\) by graded \(\Sym(\underline{E})\)-Modules \(\underline{L}_i\) locally free of finite type: \(0\longrightarrow \underline{L}_{r+1}\longrightarrow\ldots\longrightarrow \underline{L}_0\longrightarrow \underline{M}\longrightarrow 0\).

\item The \(\underline{L}_i\) admit a finite filtration whose quotients are of the form \(\underline{N}_j\otimes_{\underline{O}_S}\Sym(\underline{E})(j)\), where the \(\underline{N}_j\) are locally free of finite type on \(S\).
\end{enumerate}

%%%%%%%%%%  scan idx 379  |  folio 366  %%%%%%%%%%
\SGAPrintPage{366}%
\(2^\circ\)) \(1,\ldots,\xi^{r+1}\) are linearly dependent over \(K^{\bullet}(S)\).

\(3^\circ\)) \(1,\ldots,\xi^r\) are linearly independent over \(K^{\bullet}(S)\).

\sgaStableTarget{sga6.e06.sec.1.2}1.2.
For convenience, we shall recall here a few classical definitions and properties of graded Modules.

Let \(S\) be a scheme, \(\underline{B}\) a graded \(\underline{O}_S\)-Algebra, \(\underline{M}\) a graded \(\underline{B}\)-Module, For every graded \(\underline{B}\)-Module \(\underline{N}\), we denote as usual by \(\underline{N}(n)\) the graded \(\underline{B}\)-Module defined by:
\[
  (\underline{N}(n))_k=\underline{N}_{n+k}\quad\text{for every }k.
\]
We say that \(\underline{M}\) is a locally free graded \(\underline{B}\)-Module (resp. locally free of finite type) if every point of \(S\) has an open neighborhood \(U\) such that \(\underline{M}|U\) is isomorphic, as a graded \(\underline{B}\)-Module, to a direct sum of Modules of the form \(\underline{B}(n)|U\) (resp. a finite direct sum).

We say that \(\underline{M}\) is a graded \(\underline{B}\)-Module of finite presentation if every point of \(S\) has an open neighborhood \(U\) such that there is an exact sequence \(\underline{P}\longrightarrow \underline{Q}\longrightarrow \underline{M}|U\longrightarrow 0\), where \(\underline{P}\) and \(\underline{Q}\) are free graded \((\underline{B}|U)\)-Modules of finite type, and the homomorphisms are compatible with the gradings.

1.2.1. For every integer \(n\), we can define a canonical homomorphism of graded \(\underline{B}\)-Modules:
\sgaStableTarget{sga6.e06.eq.1-2-1}\begin{equation}\tag{1.2.1}
  \underline{M}_n\otimes_{\underline{O}_S} \underline{B}(-n)\longrightarrow \underline{M},
\end{equation}
by sending \(x\otimes a\), where \(x\) is a section of \(\underline{M}_n\), and \(a\) a section of \((\underline{B}(-n))_k=\underline{B}_{k-n}\), to the element \(a.x\) of \(\underline{M}_k\).

1.2.2. We say that a graded \(\underline{B}\)-Module \(\underline{M}\) is free of finite type and homogeneous of degree \(n\) if it is isomorphic as a graded \(\underline{B}\)-Module to \((\underline{B}(-n))^k\) for some integer \(k\).

%%%%%%%%%%  scan idx 380  |  folio 367  %%%%%%%%%%
\SGAPrintPage{367}%
Let \(\underline{M}\) be a free graded \(\underline{B}\)-Module of finite type. By 1.2, we can find a graded isomorphism:
\sgaStableTarget{sga6.e06.eq.1-2-2}\begin{equation}\tag{1.2.2}
  \underline{M}\simeq \bigoplus_{n=n_0}^{n_1} \bigl(\underline{B}(-n)\bigr)^{k_n}.
\end{equation}

We thus obtain a decomposition of \(\underline{M}\) as a direct sum of free graded \(\underline{B}\)-Modules of finite type that are homogeneous, namely:
\[
  \underline{M}\simeq \bigoplus_{n=n_0}^{n_1}\underline{M}^{(n)},
\]
where \(\underline{M}^{(n)}\) is free homogeneous of degree \(n\) and of rank \(k_n\).

\sgaStableTarget{sga6.e06.prop.1-2}\textbf{Proposition 1.2.3.} \uline{The integers} \(k_n\) \uline{occurring in decomposition} (1.2.2) \uline{do not depend on the chosen decomposition}.

Suppose given two decompositions, hence a graded isomorphism:
\[
  \bigoplus_{n=n_0}^{n_1}\bigl(\underline{B}(-n)\bigr)^{k_n}\xrightarrow{\sim}
  \bigoplus_{n=n'_0}^{n'_1}\bigl(\underline{B}(-n)\bigr)^{k'_n}.
\]
The proof proceeds by induction on the number of nonzero free homogeneous direct summands occurring in the decomposition on the left. If this is \(0\), then \(\underline{M}=0\), and the proposition is evident. Suppose it proved for \(n_1-n_0=m\), and let us prove it for \(n_1-n_0=m+1\). The given isomorphism, being compatible with the gradings, gives an isomorphism on the homogeneous components of the two sides, and in particular an isomorphism of the homogeneous components of degree \(<n_0\), which shows that for \(n<n_0\), we have \(k'_n=0\), and on the homogeneous components of degree \(n_0\), which shows that they have the same rank, and consequently that \(k_{n_0}=k'_{n_0}\). The given isomorphism therefore induces an isomorphism of the free homogeneous direct summands of lowest degree in the two decompositions, and hence passes to the quotient by these direct summands; we then apply the induction hypothesis to the quotient Module, whence the result.

%%%%%%%%%%  scan idx 381  |  folio 368  %%%%%%%%%%
\SGAPrintPage{368}%
It follows in particular from 1.2.3 that for every free graded Module of finite type \(\underline{M}\), the number of elements of a homogeneous basis of \(\underline{M}\) having a given degree does not depend on the chosen basis. In particular, the integers \(n_0\) and \(n_1\) occurring in decomposition (1.2.2), defined by the conditions \(k_{n_0}\neq0\), \(k_{n_1}\neq0\), do not depend on the decomposition.

\sgaStableTarget{sga6.e06.prop.1-3}\textbf{Proposition 1.3.} \uline{Let} \(S\) \uline{be a quasi-compact scheme}, \(\underline{E}\) \uline{a} \(\underline{O}_S\)\uline{-Module locally free of finite type}, \(X\) \uline{the associated projective bundle}. \uline{Let} \(\underline{F}\) \uline{be a} \(\underline{O}_X\)\uline{-Module of finite presentation}, \uline{flat over} \(S\). \uline{There exists an integer} \(n_0\), \uline{such that for} \(n\geq n_0\), \uline{we have}:
\[
\begin{array}{ll}
\text{i)}& R^i f_*(\underline{F}(n))=0\quad\text{\uline{for} }i\geq 1;\\[2mm]
\text{ii)}& f_*(\underline{F}(n))\text{ \uline{is locally free of finite type on} }S;\\[2mm]
\text{iii)}& \underline{M}=\displaystyle\bigoplus_{n\geq n_0} f_*(\underline{F}(n))\text{ \uline{is a graded} }\Sym(\underline{E})\text{\uline{-Module of finite presentation} (1.2).}
\end{array}
\]

The desired properties are local on \(S\); since \(S\) is quasi-compact, we can, by taking the largest of the integers found, reduce to the case where \(S\) is affine.

First suppose that \(S\) is noetherian.

Property i) then follows from Serre's theorem.

Property ii) follows from EGA III 7.9.10, with \(n_0\) chosen so that property i) holds.

By EGA II 3.4.5, there exists a quasi-coherent graded \(\Sym(\underline{E})\)-Module \(\underline{M}'\), of finite type, such that \(\underline{F}=\Projz(\underline{M}')\). Moreover, \(\Gamma_*(\underline{F})=\Gamma_*(\Projz \underline{M}')=\displaystyle\bigoplus_{n\in\mathbb Z} f_*(\underline{F}(n))\) is TN-isomorphic to \(\underline{M}'\) by EGA III 2.3.1; thus there exists an integer \(n_1\) such that \(\displaystyle\bigoplus_{n\geq n_1} f_*(\underline{F}(n))\simeq \bigoplus_{n\geq n_1} \underline{M}'_n\),
%%%%%%%%%%  scan idx 382  |  folio 369  %%%%%%%%%%
\SGAPrintPage{369}%
and hence such that \(\bigoplus_{n\geq n_1} f_*(\underline{F}(n))\) is a \(\Sym(\underline{E})\)-Module of finite type, and therefore of finite presentation by the noetherian hypotheses; the result moreover holds for every integer \(n_1\), since the \(f_*(\underline{F}(n))\) are coherent by Serre's theorem.

The proposition is therefore proved in the case where \(S\) is noetherian.

Suppose now that we have a Cartesian square:
\[
\begin{tikzcd}[column sep=large,row sep=large]
X \arrow[r,"g"] \arrow[d,"f"'] & X_0 \arrow[d,"f_0"]\\
S \arrow[r,"g_0"] & S_0,
\end{tikzcd}
\]
with \(S_0\) noetherian, \(X_0=\mathbf P(\underline{E}_0)\); let \(\underline{F}_0\) be a \(\underline{O}_{X_0}\)-Module of finite presentation, flat over \(S_0\), and let \(\underline{F}=g^*(\underline{F}_0)\). By the preceding argument the proposition holds for \(\underline{F}_0\). By EGA III 7.7.5 and 7.8.5, the base-change theorem holds under the preceding hypotheses; we immediately deduce that the theorem holds for \(\underline{F}\).

To complete the proof, it therefore suffices to show that the general case (with \(S\) affine) can be obtained from the noetherian case by base change.

Since \(S\) is affine, \(S\) is a filtered projective limit of noetherian affine schemes (if \(S=\Spec(A)\), we write \(A\) as the inductive limit of its subalgebras of finite type over \(\mathbb Z\)). Thus let \(S=\varprojlim S_\lambda\). By EGA IV 8.5.2 and 8.5.5, there exists an index \(\lambda\) and a \(\underline{O}_{S_\lambda}\)-Module locally free of finite type \(\underline{E}_\lambda\) such that \(\underline{E}=u_\lambda^*(\underline{E}_\lambda)\) (denoting by \(u_\lambda\) the morphism \(S\to S_\lambda\), and by \(u_{\lambda\mu}\) the transition morphisms \(S_\mu\to S_\lambda\)).

If we put, for \(\mu\geq\lambda\), \(\underline{E}_\mu=u_{\lambda\mu}^*(\underline{E}_\lambda)\), and \(X_\mu=\mathbf P(\underline{E}_\mu)\), then for every \(\mu\), we have \(X=X_\mu\times_{S_\mu}S\), and \(X=\varprojlim X_\mu\). Consequently, there exists an index \(\lambda'\geq\lambda\)
